\documentclass[11pt]{article}
\usepackage{paperreview}
\title{Which Collision Laws Admit Finite-Moment Prediction?\\
A Gaussian-Mixture Classification and a Clock-Free Inverse}
\author{Luqman Mahmood}
\date{1 October 2026}
\hypersetup{pdftitle={Which Collision Laws Admit Finite-Moment Prediction? A Gaussian-Mixture Classification and a Clock-Free Inverse},pdfauthor={Luqman Mahmood}}
\begin{document}
\maketitle
\draftstatus{SM3, v1.0.0}{Conditional classification with exact controls; no physical-gas validation.}
\begin{abstract}
For homogeneous isotropic elastic collisions with uniform outgoing relative
direction, the initial fourth velocity-moment response of a centered
Gaussian mixture is a quadratic functional of its temperature law.
We characterize when that functional is exactly determined by any finite
list of temperature moments, allowing an arbitrary determining function.
For moments through order $m\geq2$, this occurs precisely when the thermal
collision response is a polynomial of degree at most $m-2$. Under stated
integrability assumptions, this is equivalent to a radial collision kernel
polynomial in squared relative speed. The statement concerns one initial
response, not autonomous closure of a full moment hierarchy.
On a known compact temperature interval, polynomial approximation also
bounds initial-response prediction error and the ambiguity between
preparations sharing the retained moments.
In a separately postulated Maxwell-type moment model, a two-time logarithmic ratio cancels an
unknown changing collision clock. It identifies a constant angular
coordinate, but only its exposure-weighted mean when the law varies.
An exact counterexample shows why finitely many endpoints cannot determine
the intervening law. No laboratory validation is involved.
\end{abstract}

\section{Scope and established baseline}
Closed Maxwell-molecule moment hierarchies are classical, with exact
solutions and modern finite-factor formulations \cite{Santos,Karlin}.
The historical comparators are the Maxwell-molecule moment calculations
of Ikenberry and Truesdell, Grad's moment method, and the exact
Boltzmann solution of Krook and Wu \cite{IkenberryTruesdell,Grad,KrookWu}.
They are baselines, not consequences of the present classification.
Polynomial moment preservation also has a broad independent literature
\cite{CKT}. Neither the existence of Maxwell closure nor the inadequacy
of an arbitrary truncation is offered as a new result here.
Entropy-based moment-closure hierarchies \cite{Levermore} concern a
different task from exact constancy of one response over all measures
with prescribed moments; this theorem does not rule out those closures
as useful approximations.
The truncated moment problem asks which measures represent prescribed
moments \cite{CurtoFialkow}; here the question is constancy of a particular
quadratic functional across all such measures. Laplace-transform uniqueness
is also established machinery \cite{SSV}, not a claimed new inverse theorem.

Our first question is more restricted and has a converse: which collision
responses can be predicted exactly from finitely many instantaneous
moments throughout the entire class of positive-temperature Gaussian
mixtures? Our second question concerns a different observation resource:
what does a two-time ratio identify if the collision rate or angular law
changes? Combining these questions makes explicit the distinction between
a state summary, a dynamical law and an observation history.

\section{Preparations and the initial response}
Let $\mu$ be a compactly supported probability measure on $(0,\infty)$ and
\[
 f_\mu(v)=\int (2\pi\theta)^{-3/2}
       \exp(-|v|^2/(2\theta))\,\mu(\dd\theta),\qquad
 M_j=\int\theta^j\,\mu(\dd\theta).
\]
Then $\int|v|^{2j}f_\mu(v)\dd v=(2j+1)!!M_j$.
Temperature moments through order $m$ correspond to radial velocity moments
through degree $2m$.

Fix the collision normalization explicitly. For a test function $\psi$,
\[
 \frac{\dd}{\dd t}\langle\psi\rangle
 =\frac12\iint f(v)f(w)K(|v-w|)
   \int_{S^2}[\psi(v')+\psi(w')-\psi(v)-\psi(w)]
              \frac{\dd\sigma}{4\pi}\,\dd v\,\dd w ,
\]
where $V=(v+w)/2$, $g=v-w$ and $v',w'=V\pm|g|\sigma/2$.
Thus $D$ below is the derivative of the raw moment
$\langle|v|^4\rangle$, not of that moment divided by $15$.
For this model the initial fourth velocity-moment derivative is
\begin{equation}\label{eq:D}
 D(\mu)=-\iint (\theta-\eta)^2\Phi(\theta+\eta)
                        \,\mu(\dd\theta)\mu(\dd\eta).
\end{equation}
For a radial speed factor $K$, the thermal response is
\begin{equation}\label{eq:transform}
 \Phi(s)=\frac{s^{-7/2}}{6\sqrt{2\pi}}
 \int_0^\infty x^{5/2}K(\sqrt{x})\e^{-x/(2s)}\dd x .
\end{equation}
These formulas assume elastic collisions and a uniform outgoing relative
direction. A nonuniform angle law is not silently absorbed into a claim
about the bare radial kernel.
For completeness, conditional on temperatures $\theta,\eta$, write
$s=\theta+\eta$. Then $g$ is Gaussian with covariance $sI$, and
$V\mid g=(\theta-\eta)g/(2s)+\zeta$, with independent centered
$\zeta$ of covariance $\theta\eta I/s$.
The angularly averaged weak increment for the fourth moment is
$K(|g|)\{ |g|^2|V|^2/3-(V\cdot g)^2\}$.
Conditioning on $g$ gives
$-(\theta-\eta)^2 K(|g|)|g|^4/(6s^2)$.
The Maxwell radial density and the substitution $x=|g|^2$ give
\eqref{eq:transform}, including its normalization.

\begin{definition}
The response is \emph{determined through order $m$} if any two admitted
probability measures sharing $M_0,\ldots,M_m$ have the same $D$.
No continuity, polynomial form or linearity of the determining map is assumed.
\end{definition}

\section{An exact finite-moment classification}
\begin{theorem}\label{thm:classification}
Let $\Phi\in C^\infty((0,\infty))$ and $m\geq2$.
On all compactly supported probability measures on $(0,\infty)$,
the following are equivalent:
\begin{enumerate}[label=(\roman*)]
\item $D$ is determined through order $m$.
\item $\Phi$ is a polynomial of degree at most $m-2$.
\end{enumerate}
For a nonzero polynomial of degree $J$, the least cutoff is $J+2$.
For a nonpolynomial smooth $\Phi$, no finite cutoff works on this class.
\end{theorem}
\begin{proof}
Write $h(x,t)=(x-t)^2\Phi(x+t)$ and
$B(\rho,\nu)=\iint h(x,t)\rho(\dd x)\nu(\dd t)$.

Suppose (i) holds. Let $\nu$ be any finite signed atomic measure with
vanishing moments through order $m$. Choose a probability measure $\rho$
giving positive mass to every support point of $\nu$. For small
$\varepsilon>0$, both $\rho+\varepsilon\nu$ and
$\rho-\varepsilon\nu$ are probability measures with the same moments.
The difference of their quadratic functionals is
$-4\varepsilon B(\rho,\nu)$, so $B(\rho,\nu)=0$.

For any $t>0$ and $0<a<1$, the measure
$\rho_a=(1-a)\rho+a\delta_t$ still permits both signed perturbations.
Thus $B(\rho_a,\nu)=0$ and consequently $B(\delta_t,\nu)=0$.
By symmetry of $h$, this says that $h(\,\cdot\,,t)$ annihilates every such
signed measure. In particular, the $(m+1)$st forward difference in $x$
vanishes, since its binomial signed measure annihilates all monomials
through order $m$. Taking its mesh to zero yields
$\partial_x^{m+1}h(x,t)=0$ for every $x,t>0$.
At $x=t$ the Leibniz rule leaves only
\[
 0=m(m+1)\Phi^{(m-1)}(2t).
\]
Connectedness of the positive half-line proves (ii).
This argument never differentiates the unknown determining map.

Conversely, for $\Phi(s)=\sum_{j=0}^J c_js^j$ direct expansion gives
\begin{equation}\label{eq:moment}
 D=-\sum_{j=0}^Jc_j\sum_{\ell=0}^j\binom j\ell
 \left[M_{\ell+2}M_{j-\ell}
       -2M_{\ell+1}M_{j-\ell+1}
       +M_\ell M_{j-\ell+2}\right].
\end{equation}
Only moments through $J+2$ occur. Necessity excludes any smaller cutoff
of at least two. A cutoff of one cannot suffice for any nonzero $\Phi$:
$\delta_a$ and $(\delta_{a-h}+\delta_{a+h})/2$ have the same mean but
responses $0$ and $-2h^2\Phi(2a)$, respectively. Choose $a$ with
$\Phi(2a)\ne0$ and $0<h<a$.
\end{proof}

The supplement \cite{Internal} checks the moment formula and separating
examples exactly; these finite checks do not replace the proof of the
universally quantified statement.

\begin{corollary}[Approximate initial-response prediction]\label{cor:approximation}
Let $0<a<b$, $m\geq2$, $\Phi\in C([2a,2b])$, and let $\mu$ be a probability measure on $[a,b]$
with temperature variance $v=M_2-M_1^2$. For a polynomial $P$ of degree
at most $m-2$, put
\[
 \varepsilon(P)=\norm{\Phi-P}_{\infty,[2a,2b]},\qquad
 D^P(\mu)=-\iint(\theta-\eta)^2P(\theta+\eta)
                         \,\mu(\dd\theta)\mu(\dd\eta).
\]
The polynomial predictor $D^P(\mu)$ is determined by $M_0,\ldots,M_m$
through \eqref{eq:moment}, and
\[
 |D(\mu)-D^P(\mu)|\leq2v\varepsilon(P).
\]
If $\nu$ is another probability measure on the same interval with those
moments, then, writing
$\varepsilon_m=\inf_{\deg P\leq m-2}\varepsilon(P)$,
\[
 |D(\mu)-D(\nu)|\leq4v\varepsilon_m
                   \leq(b-a)^2\varepsilon_m.
\]
\end{corollary}
\begin{proof}
The polynomial moment formula gives $D^P(\mu)=D^P(\nu)$, while
\[
 \iint(\theta-\eta)^2\,\mu(\dd\theta)\mu(\dd\eta)=2v.
\]
Bounding the integrand error gives the single-predictor inequality.
The triangle inequality and the common variance give the two-preparation
bound after taking the infimum over $P$. Finally
$\int(\theta-a)(b-\theta)\,\mu(\dd\theta)\geq0$ implies
$v\leq(M_1-a)(b-M_1)\leq(b-a)^2/4$.
\end{proof}
Thus $2v\varepsilon(P)$ controls a specified predictor, whereas
$4v\varepsilon_m$ bounds the diameter of responses compatible with the
retained moments. This elementary approximation consequence does not
claim that the diameter bound is sharp. A numerical guarantee requires
an actual polynomial and a certified uniform error bound on the stated
support interval, not only the definition of $\varepsilon_m$.

\subsection{What this theorem does not close}
Even when \eqref{eq:moment} holds, the higher moments appearing in its
right-hand side need not themselves have autonomous equations at that cutoff.
The theorem addresses the initial derivative on a preparation class, not the
invariance of that class under every subsequent collision evolution.
It does not exclude approximate closure, restricted-family closure,
extra measured observables or history-dependent inference.
Corollary~\ref{cor:approximation} concerns only the initial response;
uncertain support, noisy moments and propagation to later times require
separate error estimates.

\section{Classification in the raw speed variable}
\begin{corollary}\label{cor:kernel}
Assume $K\geq0$ is measurable and
$\int_0^\infty r^6K(r)\e^{-\lambda r^2}\dd r<\infty$ for every
$\lambda>0$. Then the thermal response in \eqref{eq:transform} is a
polynomial $\sum_{j=0}^Jc_js^j$ if and only if, almost everywhere on $r>0$,
\[
 K(r)=\sum_{j=0}^J\frac{6c_j}{(2j+5)!!}r^{2j}.
\]
The resulting polynomial must be nonnegative on the positive half-line;
its individual coefficients need not be nonnegative.
\end{corollary}
\begin{proof}
The Gamma integral maps $K(r)=r^{2j}$ to
$\Phi(s)=(2j+5)!!s^j/6$. Subtract this finite polynomial from $K$
when $\Phi$ is polynomial. The Laplace transform of
$x^{5/2}$ times the difference at $\sqrt{x}$ vanishes for every
positive transform parameter. The stated exponential integrability allows
Laplace uniqueness, so the difference vanishes almost everywhere.
\end{proof}
The value at $r=0$ is not identified. This is not a classification of
anisotropic scattering or of all intermolecular force laws.

\section{A clock-canceling inverse in a Maxwell model}
Consider a separate admitted angular family with normalized moments
$q_j=\langle|v|^{2j}\rangle/[(2j+1)!!\,T^j]$, where
$T=\langle|v|^2\rangle/3$ is conserved, and equations
\begin{align}
 q_2'&=-\nu(t)(q_2-1)/2,\nonumber\\
 q_3'&=-3\nu(t)(q_3-q_2)/4,\label{eq:maxwell}\\
 q_4'&=\nu(t)\left[-(6+G)q_4/8+Gq_3/2
                              +3(2-G)q_2^2/8\right].\nonumber
\end{align}
Here $\nu\geq0$ and $\nu G$ are integrable on each observation interval,
and $G(t)$ is an admitted angular coordinate.
These equations, rather than an unqualified identification of $G$ with
a molecular parameter, define the inverse problem.

Set $y=q_2-1$ and $R=q_4-4q_3+3q_2^2$.
Substitution gives
\[
 y'=-\nu y/2,\qquad R'=-\nu(6+G)R/8.
\]
\begin{theorem}\label{thm:clock}
If $y,R>0$ and $A=\int_{t_1}^{t_2}\nu(t)\dd t>0$, then
\begin{equation}\label{eq:inverse}
 \widehat G:=
 4\frac{\log(R(t_2)/R(t_1))}{\log(y(t_2)/y(t_1))}-6
 =\frac{\int_{t_1}^{t_2}\nu(t)G(t)\dd t}{A}.
\end{equation}
A constant $G$ is recovered despite an unknown changing common clock.
For varying $G$, only the collision-exposure-weighted mean is identified.
\end{theorem}
\begin{proof}
Integrate the two scalar logarithmic derivatives and eliminate $A$.
\end{proof}
The denominator and sign conditions are admission requirements.
For zero exposure or a vanishing invariant the displayed inverse is
undefined. If uncertainty intervals permit either case or an invalid
logarithm, a central estimate alone cannot justify applying the inverse.

\subsection{Exact endpoint ambiguity}
Use the same initial $q_2,q_3,q_4$ and total exposure in two experiments.
In the first take $G=7/5$ throughout. In the second take $G=6/5$ for the
first half of exposure and $G=8/5$ for the second.
Both histories have the same exposure-weighted mean. The first two moment
equations are independent of $G$; the equation for $R$ depends only on
its integrated value. Therefore every displayed endpoint moment is identical.
An intermediate observation can distinguish these particular histories,
but the same construction can be placed inside each interval of any
finite schedule. This defeats unrestricted history identification, not
constant-law inference.

Exact continuous data instead supply the exposure coordinate
$s=-2\log(y/y(t_1))$, along which
\[
 \frac{\dd\log R}{\dd s}=-(6+G(s))/8.
\]
Recovering this derivative from noisy finite data is a separate stability
problem, not solved by \eqref{eq:inverse}.

\section{Predictive interpretation and a finite next test}
A useful future experiment would first prepare two independently verified
velocity distributions with the same admitted low moments. A specified
ordinary collision law would predict their unequal initial fourth-moment
responses before either slope was observed. A separate rate calibration,
or the admitted two-time inverse, would constrain the changing clock.

The classification says when equality of finitely many summaries can
guarantee equality of the initial response. It does not say how large a
detectable separation is for every preparation or noise budget.

\section{Reproducibility and remaining review}
The supplement supplies an exact-rational counterexample with matching
$M_0,M_1,M_2$, polynomial-response controls, the transform constants,
and a symbolic check of the clock-canceling invariant and endpoint
ambiguity. No private data or machine-specific path is required.
Its finite checks complement the proofs; they do not establish the
converse by enumeration. The cited works establish the relevant baselines;
a systematic novelty comparison for this particular classification has
not been made, and no priority is claimed.

\aistatement
\begin{thebibliography}{99}
\bibitem{IkenberryTruesdell} E. Ikenberry and C. Truesdell,
\emph{On the pressures and the flux of energy in a gas according to
Maxwell's kinetic theory, I}, J. Rational Mech. Anal. \textbf{5}, 1--54 (1956).
\url{https://doi.org/10.1512/iumj.1956.5.55001}.
\bibitem{Grad} H. Grad,
\emph{On the kinetic theory of rarefied gases},
Comm. Pure Appl. Math. \textbf{2}, 331--407 (1949).
\url{https://doi.org/10.1002/cpa.3160020403}.
\bibitem{KrookWu} M. Krook and T. T. Wu,
\emph{Formation of Maxwellian Tails},
Phys. Rev. Lett. \textbf{36}, 1107--1109 (1976).
\url{https://doi.org/10.1103/PhysRevLett.36.1107}.
\bibitem{Santos} A. Santos,
\emph{Solutions of the moment hierarchy in the kinetic theory of Maxwell models},
Continuum Mech. Thermodyn. \textbf{21}, 361--387 (2009).
\url{https://arxiv.org/abs/0903.1054}.
\bibitem{Karlin} I. Karlin,
\emph{Fock-Space Formulation of the Boltzmann Collision Operator for Maxwell
Molecules} (2026). \url{https://arxiv.org/abs/2609.00040}.
\bibitem{CKT} C. Cuchiero, M. Keller-Ressel and J. Teichmann,
\emph{Polynomial processes and their applications to mathematical finance},
Finance Stoch. \textbf{16}, 711--740 (2012).
\url{https://doi.org/10.1007/s00780-012-0188-x}.
\bibitem{CurtoFialkow} R. E. Curto and L. A. Fialkow,
\emph{Truncated K-moment problems in several variables},
J. Operator Theory \textbf{54}, 189--226 (2005).
\url{https://arxiv.org/abs/math/0507067}.
\bibitem{SSV} R. L. Schilling, R. Song and Z. Vondra\v{c}ek,
\emph{Bernstein Functions: Theory and Applications}, second ed.,
De Gruyter (2012). \url{https://doi.org/10.1515/9783110269338}.
\bibitem{Levermore} C. D. Levermore,
\emph{Moment closure hierarchies for kinetic theories},
J. Stat. Phys. \textbf{83}, 1021--1065 (1996).
\url{https://doi.org/10.1007/BF02179552}.
\bibitem{Internal} L. Mahmood, Reproducibility supplement to this paper (2026),
\texttt{supplement/README.md} and \texttt{check.py}.
\end{thebibliography}
\end{document}
